%% FullCourse_10Lecs -- Lecture 8, Topic 8.4: Augmented Generation + RAG vs Fine-Tune + Apps + Failures + Bridge
%% Source: ./archive_pre_split/Lec08_RAG.tex (sections 9-13)
%% Pass B (2026-05-17): diffed against MiniCourse M4_T2 (RAG-vs-FT, apps, eval).
%%   FullCourse T4 is already a superset (5 use cases + RAGAS + open challenges
%%   + agentic bridge). No new frames ported; deck is at the 30-page ceiling.

%% Shared preamble for the FullCourse_10Lecs topic decks.
%%
%% Each topic .tex \input's this file, then defines its own \title, \date
%% override (if any), \graphicspath, and \begin{document}.
%%
%% Reconstructed verbatim from the inlined copy preserved in
%% Lec10_Case_Studies/Lec10_T8_Case6_DeepRamsey.tex (lines 23-190), which is
%% explicitly annotated there as "from ../shared_preamble.tex, copied verbatim".
%% Base preamble matches MiniCourse_8hr/Slides/shared_preamble.tex; this version
%% additionally provides \sectiondivider, the conceptbox/cautionbox/examplebox/
%% takeawaybox tcolorboxes, and \compacttables used across the split decks.

\documentclass[10pt,english,aspectratio=169]{beamer}

\usepackage{ae,aecompl}
\usepackage[T1]{fontenc}
\usepackage{textcomp}
\usepackage[utf8]{inputenc}
\usepackage{babel}
\usepackage{datetime2}

\setcounter{secnumdepth}{3}
\setcounter{tocdepth}{3}

\usepackage{amsmath, amssymb, amsthm, graphicx}
\usepackage{booktabs, multirow}
\usepackage{tikz}
\usepackage{pgfplots}
\pgfplotsset{compat=1.16}
\usepackage[most]{tcolorbox}
\tcbuselibrary{skins, breakable, theorems}
\usepackage{listings}
\usepackage{xcolor}

\lstset{
  basicstyle=\ttfamily\small,
  keywordstyle=\color{blue!80!black}\bfseries,
  commentstyle=\color{green!50!black}\itshape,
  stringstyle=\color{orange!80!black},
  backgroundcolor=\color{gray!8},
  frame=single,
  rulecolor=\color{gray!40},
  breaklines=true,
  showstringspaces=false,
  numbers=none,
}

\lstdefinestyle{pythonstyle}{
    language=Python,
    basicstyle=\ttfamily\footnotesize,
    keywordstyle=\color{blue!80!black}\bfseries,
    commentstyle=\color{green!50!black}\itshape,
    stringstyle=\color{orange!80!black},
    frame=single,
    breaklines=true,
    showstringspaces=false,
    numbers=left,
    numbersep=5pt,
    numberstyle=\tiny\color{gray},
    backgroundcolor=\color{gray!8},
    rulecolor=\color{gray!40}
}

\ifx\hypersetup\undefined
  \AtBeginDocument{%
    \hypersetup{bookmarks=false,breaklinks=true,pdfborder={0 0 0},colorlinks=true,
      linkcolor=DarkText,urlcolor=RedTitle}%
  }
\else
  \hypersetup{bookmarks=false,breaklinks=true,pdfborder={0 0 0},colorlinks=true,
    linkcolor=DarkText,urlcolor=RedTitle}%
\fi

\makeatletter
\newcommand\makebeamertitle{%
  \begin{frame}[plain]
    \vspace*{1.0cm}
    \centering
    {\usebeamerfont{title}\usebeamercolor[fg]{title}\inserttitle\par}
    \vspace{1.0cm}
    {\usebeamerfont{author}\usebeamercolor[fg]{author}\insertauthor\par}
    \vspace{0.2cm}
    {\usebeamerfont{date}\usebeamercolor[fg]{date}\insertdate\par}
  \end{frame}}%
\AtBeginDocument{%
  \let\origtableofcontents=\tableofcontents
  \def\tableofcontents{\@ifnextchar[{\origtableofcontents}{\gobbletableofcontents}}
  \def\gobbletableofcontents#1{\origtableofcontents}%
}

\usetheme{metropolis}
\usefonttheme{professionalfonts}

\definecolor{LightGray}{RGB}{230,230,230}
\definecolor{RedTitle}{RGB}{0,0,200}
\definecolor{VeryLightGray}{RGB}{245,245,245}
\definecolor{DarkText}{RGB}{0,0,0}

\setbeamercolor{background canvas}{bg=white}
\setbeamercolor{title}{fg=RedTitle}
\setbeamercolor{author}{fg=DarkText}
\setbeamercolor{date}{fg=DarkText}
\setbeamercolor{frametitle}{bg=VeryLightGray,fg=DarkText}
\setbeamercolor{normal text}{fg=DarkText}
\setbeamerfont{title}{size=\large}

\setbeamersize{text margin left=5mm, text margin right=5mm}
\addtobeamertemplate{frametitle}{}{\vspace{-0.5em}}
\newcommand{\reducevspace}{\vspace{-0.5cm}}

% Shared section divider and box styles for split topic decks.
\newcommand{\sectiondivider}[2]{%
\begin{frame}[plain,noframenumbering]
\begin{center}
\vfill
{\Large\textbf{#1}\par}
\vspace{0.35cm}
{\normalsize #2\par}
\vfill
\end{center}
\end{frame}
}

\newtcolorbox{conceptbox}[1][]{
  colback=blue!3,
  colframe=RedTitle!55,
  boxrule=0.35mm,
  arc=1mm,
  left=2mm,
  right=2mm,
  top=1mm,
  bottom=1mm,
  #1
}
\newtcolorbox{cautionbox}[1][]{
  colback=red!3,
  colframe=red!50!black,
  boxrule=0.35mm,
  arc=1mm,
  left=2mm,
  right=2mm,
  top=1mm,
  bottom=1mm,
  #1
}
\newtcolorbox{examplebox}[1][]{
  colback=green!3,
  colframe=green!45!black,
  boxrule=0.35mm,
  arc=1mm,
  left=2mm,
  right=2mm,
  top=1mm,
  bottom=1mm,
  #1
}
\newtcolorbox{takeawaybox}[1][]{
  colback=VeryLightGray,
  colframe=RedTitle!55,
  boxrule=0.35mm,
  arc=1mm,
  left=2mm,
  right=2mm,
  top=1mm,
  bottom=1mm,
  #1
}
\newcommand{\compacttables}{\renewcommand{\arraystretch}{1.15}}

\setbeamertemplate{navigation symbols}{}
\setbeamertemplate{footline}{%
  \hfill%
  \usebeamercolor[fg]{page number in head/foot}%
  \usebeamerfont{page number in head/foot}%
  \insertframenumber\,/\,\inserttotalframenumber\kern1em\vskip2pt%
}

\makeatother

\author{Zhigang Feng}
% "date with time": compile timestamp so each build is identifiable.
\date{\today\ \DTMcurrenttime}


\usetikzlibrary{arrows.meta, positioning, shapes.geometric, fit, calc, backgrounds}

\graphicspath{{./pic/}{./figures/}{../../MiniCourse_8hr/Slides/Module4_Agentic_AI/pic/}{../}}

\title[AI for Econ Research]{AI for Economic Research: Dynamic Models, Language, and Agents\\
Lecture 8 --- Topic 8.4: Augmented Generation, Applications, Failure Modes, and Bridge}

\begin{document}

\makebeamertitle

%=============================================================================
\begin{frame}{Topic 8.4 Agenda}
\begin{enumerate}
  \item \textbf{Augmented Generation}
  \medskip
  \item \textbf{RAG vs Fine-Tuning}
  \medskip
  \item \textbf{Economic Applications}
  \medskip
  \item \textbf{OLG RAG Demo}
  \medskip
  \item \textbf{Failure Modes and Evaluation}
  \medskip
  \item \textbf{Bridge and Summary}
\end{enumerate}
\end{frame}

%=============================================================================
\section{Augmented Generation}

\begin{frame}{From Retrieved Chunks to a Final Answer}
\begin{itemize}
\item We now have the top-$K$ chunks for the user's query (typically $K = 3$ to $10$). What do we do with them?

\medskip
\item \textbf{Insert them into a prompt template, send to the LLM.} The model sees:
\begin{enumerate}
    \item System instructions (how to behave)
    \item The retrieved chunks (the evidence)
    \item The user question
    \item Output-format instructions (how to cite)
\end{enumerate}

\medskip
\item The LLM now \textbf{generates} an answer that is \emph{grounded in the retrieved evidence}.

\medskip
\item This is the only place a generative model is used in the whole pipeline. Everything before this --- chunking, embedding, retrieval --- exists to put high-quality evidence in front of the model.

\medskip
\item \textbf{The single most important design choice} at this stage: the \emph{prompt template}.
\end{itemize}
\end{frame}

\begin{frame}[fragile]{Prompt Template --- Anatomy}
A real RAG prompt has four parts. Lines below are color-coded by role:
\begin{lstlisting}
SYSTEM:
You are a research assistant for academic economists.
Answer the user's question using ONLY the context below.
If the context does not contain the answer, reply
exactly: "I don't know based on the provided sources."
Cite sources inline using [doc_id:page] notation.

CONTEXT:
[1] (fomc_2024_06:p3)
"Participants observed that inflation remained elevated..."

[2] (fomc_2024_06:p4)
"Several noted that core services prices, especially shelter,
continued to rise..."

[3] (powell_2024_press:p2)
"The Committee remains attentive to the risks that inflation poses..."

QUESTION:
How did the FOMC discuss inflation expectations in June 2024?

INSTRUCTIONS:
Answer in 2-3 sentences. Cite each claim inline as [doc_id:page].
\end{lstlisting}
\end{frame}

\begin{frame}{Source Attribution Mechanics}
\begin{itemize}
\item \textbf{Why economists need this:} Every claim in an academic paper must be \emph{traceable to a source}. An ungrounded LLM output is not citable. A grounded output is.

\medskip
\item \textbf{How to enforce:}
\begin{enumerate}
    \item At chunking time, give every chunk a stable \texttt{doc\_id} and \texttt{page} (or \texttt{turn}, or \texttt{section}).
    \item At retrieval time, return the metadata along with the chunk text.
    \item In the prompt template, label each chunk \texttt{[1] (doc\_id:page)} and instruct the model to cite using that label.
    \item In post-processing, replace the model's inline citations with hyperlinks back to the original document.
\end{enumerate}

\medskip
\item \textbf{Before / after example:}
\begin{itemize}
    \item \emph{Ungrounded:} ``The FOMC was concerned about inflation persistence in 2024.''
    \item \emph{Grounded:} ``The FOMC was concerned about inflation persistence in mid-2024 \texttt{[fomc\_2024\_06:p3]}, particularly in core services prices including shelter \texttt{[fomc\_2024\_06:p4]}.''
\end{itemize}
\end{itemize}
\end{frame}

\begin{frame}{Refusal and Faithfulness}
\begin{itemize}
\item \textbf{``I don't know'' is a feature, not a bug.}

\medskip
\item In an academic context, \emph{fabricating an answer is far worse than refusing to answer.}

\medskip
\item \textbf{Faithfulness} (the answer is fully supported by the retrieved context) matters more than \emph{fluency} (the answer reads well). A fluent fabrication is the worst possible outcome for research.

\medskip
\item \textbf{How to enforce refusal:}
\begin{itemize}
    \item Explicit instruction in the system prompt: ``If the context does not contain the answer, reply exactly: I don't know based on the provided sources.''
    \item Set generation \texttt{temperature} low (e.g.\ 0.0--0.2) to discourage creative inference.
    \item In post-processing, run a faithfulness check (Section X) and flag answers whose claims aren't supported by retrieved chunks.
\end{itemize}

\medskip
\item \textbf{Example refusal (good behavior):}
\begin{quote}
\small\itshape
``I don't know based on the provided sources. The retrieved minutes from June 2024 do not discuss the impact of fiscal policy.''
\end{quote}

\medskip
\item This protects the researcher from quoting fiction.
\end{itemize}
\end{frame}

\begin{frame}{Worked Example: Full Pipeline Run}
\textbf{Query.} \emph{``How did Powell discuss inflation expectations in 2023 press conferences?''}

\medskip
\textbf{Pipeline trace:}
\begin{enumerate}
\item \textbf{Embed query} with \texttt{bge-large-en-v1.5} $\to$ 1024-d vector.
\medskip
\item \textbf{Retrieve top-50} via Chroma + BM25 hybrid. Filter: \texttt{doc\_type = press\_conference} AND \texttt{year = 2023}.
\medskip
\item \textbf{Re-rank} the 50 with \texttt{bge-reranker-large} $\to$ top-5.
\medskip
\item \textbf{Inject top-5 chunks} into the prompt template (previous slide), tagged \texttt{[powell\_2023\_03:p4]}, \texttt{[powell\_2023\_06:p2]}, \dots
\medskip
\item \textbf{Send to GPT-4o} at \texttt{temperature=0.1}. Receives:
\begin{quote}
\small\itshape
``In 2023 press conferences, Chair Powell repeatedly emphasized that long-run inflation expectations remained well anchored \texttt{[powell\_2023\_03:p4]}, while noting that short-run expectations had risen with realized inflation \texttt{[powell\_2023\_06:p2]}. By the end of the year, he highlighted that survey- and market-based measures had begun to ease \texttt{[powell\_2023\_12:p3]}.''
\end{quote}
\medskip
\item \textbf{Post-process:} convert inline citations to clickable links back to the original transcript pages.
\end{enumerate}
\end{frame}

%=============================================================================
\section{RAG vs Fine-Tuning}

\begin{frame}{Decision Matrix: RAG vs.\ Fine-Tuning}
\begin{center}
\renewcommand{\arraystretch}{1.3}
\small
\begin{tabular}{|p{4.6cm}|c|c|}
\hline
\textbf{Dimension} & \textbf{RAG} & \textbf{Fine-tuning} \\
\hline
Knowledge freshness (update daily) & \checkmark\checkmark & \texttimes \\
\hline
Source attribution / citations & \checkmark\checkmark & \texttimes \\
\hline
Reduces hallucination & \checkmark\checkmark & limited \\
\hline
Cost to add new documents & near zero & re-train \\
\hline
Latency at query time & moderate & low \\
\hline
Style adaptation (e.g.\ ``write like the Fed'') & \texttimes & \checkmark\checkmark \\
\hline
Output-format adaptation (structured JSON) & limited & \checkmark\checkmark \\
\hline
\end{tabular}
\end{center}

\bigskip
\textbf{Reading the table:} RAG dominates on the four dimensions that matter most for research --- freshness, attribution, hallucination control, and update cost. Fine-tuning dominates on style and format, which matter for production deployments but rarely for research.
\end{frame}

\begin{frame}{When Fine-Tuning Still Makes Sense}
\begin{itemize}
\item \textbf{Style transfer.} You want the model to \emph{write like a Federal Reserve economist} --- the cadence, the hedging, the preferred phrasing. Hand the model a thousand examples of FOMC minutes and fine-tune. RAG can't teach style.

\medskip
\item \textbf{Structured-output formatting.} You want the model to always emit a JSON dictionary of DSGE parameters. Fine-tune on $\sim$500 input/output pairs. Cheaper than prompt-engineering an enforced schema.

\medskip
\item \textbf{Specialized vocabulary.} The model needs to recognize that \texttt{LSAP} = Large-Scale Asset Purchase or that \texttt{NIPA} = National Income and Product Accounts. Some tokens may not even be in the base tokenizer; fine-tuning can teach the model to handle them gracefully.

\medskip
\item \textbf{Latency-critical settings.} Fine-tuned models embed knowledge in weights, so they don't need a retrieval round-trip. Useful for real-time applications --- but rarely a binding constraint for research.

\medskip
\item \textbf{Crucial caveat:} Fine-tuning $\neq$ teaching new \emph{facts}. Models often memorize facts unreliably during fine-tuning and still hallucinate. \emph{If you want the model to know new facts, use RAG.}
\end{itemize}
\end{frame}

\begin{frame}{The Hybrid Pattern}
\begin{itemize}
\item In production, the modern pattern is \textbf{both, in different roles}:

\medskip
\item \textbf{Fine-tune the embedding model} on a domain corpus (say, 100K FOMC sentences with synthetic query-document pairs). This is \emph{cheap} ($\sim$\$50 of compute) and improves retrieval quality dramatically.

\medskip
\item \textbf{Use a general-purpose LLM as the generator} (Claude, GPT-4o). No fine-tuning of the big model.

\medskip
\item \textbf{Use RAG as the glue.} The fine-tuned embedder ensures that the retriever finds the right chunks; the general LLM ensures fluent, faithful generation.

\medskip
\item \textbf{Why this works.} You get domain understanding (from fine-tuning) and freshness + attribution (from RAG), without the cost or fragility of fine-tuning a 70B-parameter generator.

\medskip
\item This is the dominant production architecture in 2025 for any LLM application that touches a specialized corpus --- legal, medical, financial, academic.
\end{itemize}
\end{frame}

\begin{frame}{Cost \& Operational Comparison}
\begin{center}
\renewcommand{\arraystretch}{1.35}
\small
\begin{tabular}{|p{4.0cm}|p{4.2cm}|p{4.5cm}|}
\hline
& \textbf{Fine-tuning a 7B model} & \textbf{RAG with general LLM} \\
\hline
Setup cost & $\sim$\$500 + 12 GPU-hours & $\sim$\$50 (embedding pass + Chroma) \\
\hline
Time to first answer & 2--3 days (data prep + train) & 1 afternoon \\
\hline
Cost per query & low (model is local) & API call: \$0.001--\$0.05 \\
\hline
Update with new docs & re-train weekly & add chunks in seconds \\
\hline
Source citations & no & yes \\
\hline
Skill needed & ML engineer & researcher with Python \\
\hline
\end{tabular}
\end{center}

\bigskip
\textbf{For a typical academic project} --- a few thousand queries against a curated corpus --- RAG is 10$\times$ to 100$\times$ cheaper end-to-end and gives \emph{better} answers. Fine-tuning is the right tool only when you genuinely need style or format control.
\end{frame}

%=============================================================================
\section{Economic Applications}

\begin{frame}{Use Case 1: FOMC Minutes \& Press Conferences}
\begin{itemize}
\item \textbf{Corpus.} The Federal Reserve publishes statements (post-meeting), minutes (3-week lag), press conference transcripts, and (5-year lag) full meeting transcripts. Available 1993--present at \texttt{federalreserve.gov}.

\medskip
\item \textbf{Index pipeline:}
\begin{enumerate}
    \item Scrape PDFs / HTML from the Fed archive.
    \item Parse into speaker turns (minutes) or Q\&A pairs (press conferences).
    \item Chunk with rich metadata: \texttt{meeting\_date}, \texttt{doc\_type}, \texttt{speaker}, \texttt{topic\_section}.
    \item Embed with \texttt{bge-large-en-v1.5}.
    \item Store in Chroma; run BM25 in parallel for hybrid retrieval.
\end{enumerate}

\medskip
\item \textbf{Example research queries:}
\begin{itemize}
    \item \emph{``How did the Committee's framing of inflation expectations evolve from 2021 to 2024?''}
    \item \emph{``Find every appearance of `transitory' in FOMC minutes between 2020 and 2022.''}
    \item \emph{``Which participants most often dissented on rate decisions, 2017--2024?''}
\end{itemize}

\medskip
\item \textbf{Reference:} Hansen, McMahon \& Prat (2018), \emph{QJE}, ``Transparency and Deliberation Within the FOMC: A Computational Linguistics Approach.''
\end{itemize}
\end{frame}

\begin{frame}{Use Case 2: NBER Working Papers}
\begin{itemize}
\item \textbf{Corpus.} $\sim$32{,}000 NBER working papers, 1973--present, available at \texttt{nber.org}. Each has metadata: title, abstract, authors, JEL codes, full PDF.

\medskip
\item \textbf{Use case: literature review automation.} Instead of typing keywords into Google Scholar and reading 50 abstracts, ask:
\begin{quote}
\small\itshape
``Find all NBER working papers from 2019 onwards that estimate the slope of the Phillips curve. Return title, authors, year, the central estimate, and the data source.''
\end{quote}

\medskip
\item \textbf{Pipeline:}
\begin{enumerate}
    \item Index titles + abstracts as one collection (cheap, fast, good for keyword screening).
    \item Index full-paper chunks as a second collection (richer evidence for follow-up questions).
    \item Hybrid retrieval over both. Filter on \texttt{year}, \texttt{JEL\_code}, \texttt{author}.
    \item Generator outputs a structured table from retrieved papers.
\end{enumerate}

\medskip
\item \textbf{Caveat:} Always verify citations against the actual NBER record. Even with RAG, models occasionally garble author lists or year fields.

\medskip
\item \textbf{This is one of the highest-value RAG applications for an active researcher.}
\end{itemize}
\end{frame}

\begin{frame}{Use Case 3: Congressional Record \& Legislative Text}
\begin{itemize}
\item \textbf{Corpus.} The Congressional Record contains every speech delivered on the House and Senate floors since 1873. Daily editions are available at \texttt{congress.gov}.

\medskip
\item \textbf{Use case: tracking how policy debates frame economic issues over time.}
\begin{itemize}
    \item ``How has the rhetoric around `inflation' shifted in House Banking Committee testimony from 2008 to 2024?''
    \item ``Which Senators most frequently cite Federal Reserve research in floor speeches?''
\end{itemize}

\medskip
\item \textbf{The structural challenge.} Congressional Record has a deep hierarchy --- Congress / session / chamber / date / speaker / bill / section. The Section IV chunking gotcha applies in full force. A naive chunker that ignores this hierarchy will produce essentially un-citable retrieval results.

\medskip
\item \textbf{Practical fix:} Use a Congress-specific parser (e.g.\ \texttt{congress-api-python}) to extract clean speaker turns with all metadata, then feed those into the chunker.

\medskip
\item \textbf{Reference:} Gentzkow, Shapiro \& Taddy (2019), \emph{Econometrica}, ``Measuring Group Differences in High-Dimensional Choices: Method and Application to Congressional Speech.''
\end{itemize}
\end{frame}

\begin{frame}{Use Case 4: SEC 10-K Filings (Item 1A Risk Factors)}
\begin{itemize}
\item \textbf{Corpus.} EDGAR hosts $\sim$5{,}000{,}000 SEC filings, including the annual 10-K reports of every public US firm. Each 10-K has a standardized ``Item 1A: Risk Factors'' section disclosing material risks.

\medskip
\item \textbf{Use case: cross-firm comparison of disclosed risks.}
\begin{itemize}
    \item \emph{``Compare how the four largest US banks discussed interest-rate risk in their 2022 vs.\ 2023 10-K filings.''}
    \item \emph{``Find all 10-K filings from energy firms in 2024 that mention transition risk from climate policy.''}
\end{itemize}

\medskip
\item \textbf{Pipeline:}
\begin{enumerate}
    \item Use \texttt{sec-edgar-downloader} to pull 10-Ks.
    \item Run a section-aware parser to extract Item 1A only.
    \item Chunk \emph{within} Item 1A, preserving \texttt{cik}, \texttt{ticker}, \texttt{filing\_date}, \texttt{fiscal\_year}.
    \item Store in Chroma. Hybrid retrieval with metadata filter on \texttt{ticker} for cross-firm queries.
\end{enumerate}

\medskip
\item \textbf{Reference:} Loughran \& McDonald (2011), \emph{JF}, ``When Is a Liability Not a Liability? Textual Analysis, Dictionaries, and 10-Ks.'' Lopez-Lira (2023), ``Risk Factors That Matter.''
\end{itemize}
\end{frame}

\begin{frame}{Use Case 5: Cross-Country Central Bank Speeches}
\begin{itemize}
\item \textbf{Corpus.} The Bank for International Settlements maintains a unified archive of central bank speeches from $\sim$60 institutions, 1997--present, at \texttt{bis.org/cbspeeches}. Other primary sources: ECB, BoE, BoJ, RBA, BoC, RBI, PBoC.

\medskip
\item \textbf{Use case: comparative monetary policy analysis.}
\begin{itemize}
    \item \emph{``How did major central banks frame the post-COVID inflation surge in 2022 speeches? Compare the Fed, ECB, BoE, and BoJ.''}
    \item \emph{``Which central banks most often cite financial-stability concerns alongside inflation, 2010--2024?''}
\end{itemize}

\medskip
\item \textbf{Pipeline notes:}
\begin{itemize}
    \item Add \texttt{country}, \texttt{institution}, \texttt{speaker\_role}, \texttt{language} as metadata.
    \item For non-English speeches, either use a multilingual embedder (\texttt{multilingual-e5-large}) or translate first.
    \item Hybrid retrieval is critical --- many speeches use country-specific institutional jargon a generic embedder won't recognize.
\end{itemize}

\medskip
\item \textbf{This use case is essentially impossible without RAG --- no human can read $\sim$10{,}000 speeches across $\sim$60 institutions per year.}
\end{itemize}
\end{frame}

\begin{frame}{Use Case 6: Combining Structured + Unstructured}
\begin{itemize}
\item RAG isn't only for unstructured text. The most powerful pattern combines a \textbf{structured database} (e.g.\ FRED time series, BEA NIPA tables) with an \textbf{unstructured text store} (Fed speeches, FOMC minutes).

\medskip
\item \textbf{Example query:}
\begin{quote}
\small\itshape
``In months when the unemployment rate fell below 4\%, what did the Fed publicly say about labor markets?''
\end{quote}

\medskip
\item \textbf{Pipeline:}
\begin{enumerate}
    \item \emph{Structured side:} SQL on FRED data. \texttt{SELECT date FROM unrate WHERE value < 4.0;}
    \item \emph{Unstructured side:} retrieve speeches and FOMC minutes whose \texttt{date} falls within those months.
    \item Pass the joined evidence to the LLM with a prompt that says ``summarize how Fed officials characterized the labor market in these specific months.''
\end{enumerate}

\medskip
\item \textbf{This pattern} --- structured SQL filter $\to$ semantic retrieval $\to$ generation --- is uniquely useful for empirical economists. It bridges what economists already know how to do (data) with what they couldn't do before (text at scale).
\end{itemize}
\end{frame}

\begin{frame}{Synthesis: Three Workflow Patterns}
\begin{enumerate}
\item \textbf{Literature review.} Index a paper repository (NBER, SSRN, RePEc). Query: ``find work on X, summarize, return citations.'' Saves days per project.

\medskip
\item \textbf{Policy-text monitoring over time.} Index an institutional corpus (FOMC, Congress, central banks). Query: ``how did the framing of Y evolve from year A to year B?'' Turns institutional text into a longitudinal research instrument.

\medskip
\item \textbf{Cross-sectional comparative analysis.} Index a corpus where the unit of analysis is a firm or country (10-Ks, central bank speeches). Query: ``compare how entity-1, entity-2, \dots\ discussed Z in year Y.'' Enables cross-firm and cross-country comparisons that were previously infeasible at scale.
\end{enumerate}

\bigskip
\textbf{Methodological punchline.} RAG turns institutional text from \emph{background reading} into a \emph{queryable research instrument}. That is the new capability. Everything else --- chunking, embeddings, vector stores --- is plumbing.
\end{frame}


%=============================================================================
\section{Classroom Demo: OLG RAG}

\begin{frame}{Demo: RAG Over an OLG Paper}
\textbf{Question for the classroom:} What did Diamond (1965) add to the OLG approach?

\vspace{0.25cm}
\begin{center}
\begin{tikzpicture}[
  font=\footnotesize,
  node distance=0.35cm and 0.45cm,
  box/.style={draw, rounded corners=2pt, minimum width=2.15cm, minimum height=0.7cm, align=center, fill=VeryLightGray},
  arr/.style={-{Latex[length=2mm]}, thick}
]
\node[box] (source) {source paper\\markdown};
\node[box, right=of source] (chunks) {line-addressable\\chunks};
\node[box, right=of chunks] (index) {local TF-IDF\\index};
\node[box, right=of index] (retrieve) {top-$K$\\evidence};
\node[box, right=of retrieve] (answer) {grounded\\answer};
\draw[arr] (source) -- (chunks);
\draw[arr] (chunks) -- (index);
\draw[arr] (index) -- (retrieve);
\draw[arr] (retrieve) -- (answer);
\end{tikzpicture}
\end{center}

\vspace{0.2cm}
\begin{itemize}
  \item Source: \texttt{../../MiniCourse\_8hr/demos/M4\_rag\_olg\_demo}
  \item Default path is offline: no API key, no network, no vector database service.
  \item Teaching point: citations are an answer contract, not decoration.
\end{itemize}
\end{frame}

\begin{frame}[fragile]{Demo Launch and Expected Artifacts}
\begin{lstlisting}[basicstyle=\ttfamily\scriptsize, numbers=none]
cd ../../MiniCourse_8hr/demos/M4_rag_olg_demo
/usr/bin/python3 run_demo.py
/usr/bin/python3 run.py ask "What did Diamond 1965 add to the OLG approach?"
PYTHONPATH=src /usr/bin/python3 -m unittest discover -s tests
\end{lstlisting}

\vspace{0.2cm}
\begin{tabular}{p{3.0cm}p{8.5cm}}
\toprule
\textbf{Artifact} & \textbf{What students should inspect} \\
\midrule
chunk file & boundaries and line references \\
retrieval report & top-$K$ evidence, scores, and anchor checks \\
extractive answer & whether every claim is supported by retrieved text \\
unsupported query & whether the system refuses rather than improvises \\
\bottomrule
\end{tabular}

\vspace{0.12cm}
{\scriptsize Files are written under the demo's \texttt{build/} folder, including \texttt{chunks.json} and a retrieval-report CSV.}
\end{frame}

\begin{frame}{Retrieval Diagnostics Before Generation}
\begin{columns}[T]
\begin{column}{0.48\textwidth}
\textbf{Good retrieval looks like}
\begin{itemize}\small
  \item top chunks mention the relevant model contribution
  \item evidence spans are narrow enough to quote-check
  \item retrieved lines answer the question before the LLM sees them
  \item a retrieval report can be saved and audited later
\end{itemize}
\end{column}
\begin{column}{0.48\textwidth}
\textbf{Red flags}
\begin{itemize}\small
  \item high scores attached to generic background text
  \item missing author/date or page/line metadata
  \item answer uses claims absent from top-$K$ chunks
  \item citations point to nearby but non-supporting text
\end{itemize}
\end{column}
\end{columns}

\vspace{0.25cm}
\textbf{Rule for students:} read the retrieval report before you read the generated answer. If the evidence is weak, generation cannot rescue the answer.
\end{frame}

\begin{frame}{Grounded Prompt and Refusal Contract}
\small
\begin{tabular}{p{2.9cm}p{8.7cm}}
\toprule
\textbf{Prompt part} & \textbf{Contract} \\
\midrule
Role & answer as a research assistant for economists \\
Context & use only retrieved chunks, with line identifiers preserved \\
Question & answer the user query, not a broader textbook question \\
Citation rule & cite every substantive claim with chunk or line metadata \\
Refusal rule & say the source does not support the answer when evidence is absent \\
\bottomrule
\end{tabular}

\vspace{0.3cm}
\begin{takeawaybox}
\textbf{Why this belongs in Lec. 8:} the model's behavior is downstream from the index. RAG is a research instrument only when retrieval, prompting, citation, and refusal are all visible.
\end{takeawaybox}
\end{frame}


%=============================================================================
\section{Failure Modes and Evaluation}

\begin{frame}{When RAG Fails}
\textbf{Six common failure modes:}
\begin{enumerate}
\item \textbf{Retriever miss.} Relevant document exists but the embedder ranks it outside top-$K$. Often caused by terminology mismatch.

\medskip
\item \textbf{Re-ranker promotes the wrong context.} Cross-encoder mis-scores; the LLM gets a confident-but-irrelevant chunk.

\medskip
\item \textbf{Generator ignores the retrieved context.} The LLM falls back on its parametric memory and answers from training data --- often confidently and incorrectly.

\medskip
\item \textbf{Chunk boundary cuts critical information.} The fact spans the boundary between two chunks; neither chunk alone contains the answer.

\medskip
\item \textbf{Outdated index.} The corpus has been updated, but the index hasn't been re-run. The retriever returns stale chunks.

\medskip
\item \textbf{Ambiguous query.} ``Inflation'' could mean CPI, PCE, expectations, or asset prices. The retriever picks one interpretation; the user wanted another.
\end{enumerate}
\end{frame}

\begin{frame}{Mitigations --- One per Failure Mode}
\begin{itemize}
\item \textbf{Retriever miss} $\Rightarrow$ \emph{query expansion} and \emph{HyDE} (hypothetical document embeddings: ask the LLM to draft a fake answer, embed \emph{that}, retrieve nearest documents). Also: hybrid sparse + dense retrieval.

\medskip
\item \textbf{Re-ranker mis-scores} $\Rightarrow$ try a stronger re-ranker (\texttt{bge-reranker-v2}); train a domain-specific re-ranker on labeled query-document pairs.

\medskip
\item \textbf{Generator ignores context} $\Rightarrow$ stricter system prompt (``Answer ONLY from the context''); lower temperature; faithfulness check in post-processing.

\medskip
\item \textbf{Chunk boundary cuts information} $\Rightarrow$ semantic chunking; larger chunk overlap; or use a parent-document retriever (retrieve a chunk, return its parent paragraph).

\medskip
\item \textbf{Outdated index} $\Rightarrow$ incremental ingestion pipeline triggered when source documents change; nightly re-embed for high-churn corpora.

\medskip
\item \textbf{Ambiguous query} $\Rightarrow$ query rewriting (ask the LLM to clarify or expand the query before retrieval); multi-query retrieval (generate $N$ paraphrases, retrieve for each, fuse).
\end{itemize}
\end{frame}

\begin{frame}{Evaluation: The RAGAS Framework}
\begin{itemize}
\item Es, James, Espinosa-Anke \& Schockaert (2023), ``RAGAS: Automated Evaluation of Retrieval Augmented Generation.'' Now the standard.

\medskip
\item \textbf{Four metrics, two for retrieval and two for generation:}
\begin{itemize}
    \item \textbf{Context precision.} Of the retrieved chunks, how many are actually relevant to the query? \emph{(Retriever quality.)}
    \medskip
    \item \textbf{Context recall.} Of the chunks that should have been retrieved, how many did we get? \emph{(Retriever coverage.)}
    \medskip
    \item \textbf{Faithfulness.} Are all the claims in the answer supported by the retrieved context? \emph{(Generator anti-hallucination.)}
    \medskip
    \item \textbf{Answer relevance.} Does the answer actually address the query? \emph{(Generator on-topic.)}
\end{itemize}

\medskip
\item RAGAS computes all four automatically using an LLM-as-judge. Cost a few cents per query.

\medskip
\item \textbf{The right way to use it:} build a small ($\sim$50-question) eval set for your specific corpus, run RAGAS each time you change a pipeline component, and only ship a change when the metrics improve.
\end{itemize}
\end{frame}

\begin{frame}{Open Challenges}
\begin{itemize}
\item \textbf{Multi-hop reasoning.} Questions that require synthesizing evidence across \emph{multiple} retrieved documents (``Did the Fed's response in 2008 differ from the ECB's in 2011?''). Current RAG often retrieves either one side or the other, not both. Active research area.

\medskip
\item \textbf{Long-document QA.} Documents that are themselves longer than the context window (a 600-page legal opinion). Hierarchical retrieval and recursive summarization help but aren't solved.

\medskip
\item \textbf{Tables and figures.} Most embedders are trained on prose. Tables and figures in PDFs are routinely missed or garbled. Specialized layout-aware extractors (\texttt{nougat}, \texttt{marker}) help.

\medskip
\item \textbf{Multilingual retrieval.} Retrieving from a corpus that mixes English, French, Spanish, Mandarin\dots requires multilingual embedders (\texttt{multilingual-e5-large}, \texttt{bge-m3}). Quality is uneven across languages.

\medskip
\item \textbf{Dynamic knowledge and graph maintenance.} A flat vector index is easy to append to. A graph-based index may require re-extracting entities, re-computing communities, and re-writing summaries each time new documents arrive.

\medskip
\item \textbf{Evaluation cost.} RAGAS-style LLM-as-judge eval is expensive at scale. Cheap, reliable retrieval metrics for domain-specific corpora remain an open problem.
\end{itemize}
\end{frame}

%=============================================================================
\section{Bridge and Summary}

\begin{frame}{Single-Shot RAG vs.\ Agentic RAG (Preview of Lec.\ 9)}
\begin{itemize}
\item \textbf{Standard RAG} (everything we covered today) is a \emph{linear pipeline}: one query in, one retrieve-then-generate cycle, one answer out.

\medskip
\item \textbf{Agentic RAG} adds a \emph{loop}: the LLM \emph{decides} when to retrieve, what to retrieve, whether the retrieved evidence is sufficient, and whether to retrieve again with a refined query.

\medskip
\item \textbf{Two pictures:}
\end{itemize}

\begin{center}
\begin{tikzpicture}[
  font=\footnotesize,
  node distance=0.4cm and 0.5cm,
  box/.style={draw, rounded corners=2pt, minimum width=1.4cm, minimum height=0.6cm, align=center, fill=VeryLightGray},
  arr/.style={-{Latex[length=2mm]}}
]
  % Linear
  \node[font=\scriptsize\bfseries] at (-0.2, 1.5) {Standard RAG (linear)};
  \node[box] (q1) at (0.2, 0.7) {Query};
  \node[box, right=of q1] (r1) {Retrieve};
  \node[box, right=of r1] (g1) {Generate};
  \node[box, right=of g1] (a1) {Answer};
  \draw[arr] (q1) -- (r1);
  \draw[arr] (r1) -- (g1);
  \draw[arr] (g1) -- (a1);

  % Agentic
  \node[font=\scriptsize\bfseries] at (8.0, 1.5) {Agentic RAG (loop)};
  \node[box] (q2) at (8.5, 0.7) {Query};
  \node[box, right=of q2] (p2) {Plan};
  \node[box, right=of p2] (r2) {Retrieve};
  \node[box, right=of r2] (g2) {Reason};
  \node[box, right=of g2] (a2) {Answer};
  \draw[arr] (q2) -- (p2);
  \draw[arr] (p2) -- (r2);
  \draw[arr] (r2) -- (g2);
  \draw[arr] (g2) -- (a2);
  \draw[arr] (g2.south) .. controls +(0,-0.7) and +(0,-0.7) .. (p2.south);
\end{tikzpicture}
\end{center}

\begin{itemize}
\item \textbf{Lec.\ 9 (terminal agents)} shows how Claude Code combines RAG with tool use for end-to-end research workflows. Today's pipeline is the building block.
\end{itemize}
\end{frame}

\begin{frame}{Summary: The Four-Step Pipeline + the Why}
\begin{enumerate}
\item \textbf{Chunk.} Split documents into self-contained, metadata-rich pieces.

\medskip
\item \textbf{Embed.} Map each chunk to a vector so semantically similar text lives nearby in $\mathbb{R}^d$.

\medskip
\item \textbf{Store \& Retrieve.} Use a vector DB (Chroma, Qdrant, FAISS) for local semantic search; add graph or hierarchical indexes when the task depends on relationships, coverage, or multi-hop reasoning.

\medskip
\item \textbf{Augmented Generation.} Insert the retrieved evidence into a prompt template that demands grounded, cited, refusable answers.
\end{enumerate}

\bigskip
\textbf{Why this matters:} A vanilla LLM hits a hard wall on freshness, scale, latency, cost, and attribution --- exactly the dimensions that matter most for academic research. RAG separates the cost of indexing from the cost of querying, giving you a system that is cheap to update, cheap to query, faithful to sources, and capable of refusing to answer.

\bigskip
\centering\emph{The LLM is the brain. The index is the memory. The retriever is the librarian.}
\end{frame}

\begin{frame}{Closing}
\begin{center}
\Large
RAG turns institutional text --- FOMC, NBER,\\
Congress, SEC, central banks --- into a\\
\textbf{queryable research instrument.}
\end{center}

\bigskip
\bigskip
\begin{itemize}
\item Today: the static pipeline. Build it once, query it for years.
\medskip
\item \textbf{Next lecture (Lec.\ 9):} how a terminal agent uses RAG as one tool among many, plans multi-step research workflows, and edits its own code along the way.
\end{itemize}

\bigskip
\centering\textbf{Questions?}
\end{frame}

\end{document}
